% from ega3-1-fr.tex
\oldpage[0\textsubscript{III}]{5}
\phantomsection
\label{chapter:0III-fr}

\begin{center}
{\large CHAPITRE 0 (suite)\footnotemark\par}
\vspace{1.5em}
{\Large\bfseries PRÉLIMINAIRES\par}
\end{center}

\footnotetext{Pour faciliter la recherche des références, on renverra
désormais aux paragraphes du chapitre 0 publiés avec le chapitre I par le
signe $0_{\mathrm I}$.}

\setcounter{section}{7}
\section{Foncteurs représentables}
\label{section:0.8-fr}

\setcounter{subsection}{0}
\subsection{Foncteurs représentables}
\label{subsection:0.8.1-fr}

\begin{env}[8.1.1]
\label{0.8.1.1-fr}
Nous désignerons par $\mathrm{Ens}$ la catégorie des ensembles. Soit $C$ une
catégorie~; pour deux objets $X$, $Y$ de $C$, nous poserons
$h_X(Y)=\operatorname{Hom}(Y,X)$~; pour tout morphisme $u:Y\to Y'$ dans $C$,
nous désignerons par $h_X(u)$ l'application $v\mapsto vu$ de
$\operatorname{Hom}(Y',X)$ dans $\operatorname{Hom}(Y,X)$. Il est immédiat
qu'avec ces définitions, $h_X:C\to\mathrm{Ens}$ est un \emph{foncteur
contravariant}, c'est-à-dire un objet de la catégorie, notée
$\operatorname{Hom}(C^0,\mathrm{Ens})$, des foncteurs covariants de la
catégorie $C^0$, duale de la catégorie $C$, dans la catégorie
$\mathrm{Ens}$ (T, 1.7, d) et [29]).
\end{env}

\begin{env}[8.1.2]
\label{0.8.1.2-fr}
Soit maintenant $w:X\to X'$ un morphisme dans $C$~; pour tout $Y\in C$ et
tout $v\in\operatorname{Hom}(Y,X)=h_X(Y)$, on a
$wv\in\operatorname{Hom}(Y,X')=h_{X'}(Y)$~; désignons par $h_w(Y)$
l'application $v\mapsto wv$ de $h_X(Y)$ dans $h_{X'}(Y)$. Il est immédiat
que pour tout morphisme $u:Y\to Y'$ dans $C$, le diagramme
\[
  \xymatrix{
    h_X(Y')\ar[r]^{h_X(u)}\ar[d]_{h_w(Y')} &
    h_X(Y)\ar[d]^{h_w(Y)}\\
    h_{X'}(Y')\ar[r]_{h_{X'}(u)} &
    h_{X'}(Y)
  }
\]
est commutatif~; autrement dit, $h_w$ est un \emph{morphisme fonctoriel}
$h_X\to h_{X'}$ (T, 1.2), ou encore un morphisme dans la catégorie
$\operatorname{Hom}(C^0,\mathrm{Ens})$ (T, 1.7, d)). Les définitions de
$h_X$ et de $h_w$ constituent donc la définition d'un \emph{foncteur
covariant canonique}
\[
  h:C\to\operatorname{Hom}(C^0,\mathrm{Ens}),\qquad X\mapsto h_X.
  \tag{8.1.2.1}
\]
\end{env}

\begin{env}[8.1.3]
\label{0.8.1.3-fr}
Soient $X$ un objet de $C$, $F$ un foncteur contravariant de $C$ dans
$\mathrm{Ens}$ (objet de $\operatorname{Hom}(C^0,\mathrm{Ens})$). Soit
$g:h_X\to F$ un \emph{morphisme fonctoriel}~: pour tout $Y\in C$,
